% =====================================================================
%  Árvore de máxima verossimilhança com o IQ-TREE 3
%  Prática 5 -- Disciplina Métodos Computacionais (UNESP)
%
%  Compilar com pdfLaTeX (Overleaf: Menu > Compiler > pdfLaTeX).
%
%  VERSÃO: por padrão gera a versão do ALUNO (sem respostas).
%  Para a versão do PROFESSOR (com gabarito), troque \professorfalse
%  por \professortrue na linha indicada abaixo, ou compile com
%     pdflatex "\def\versaoprofessor{}\input{apostila_iqtree.tex}"
% =====================================================================
\documentclass[11pt,a4paper]{article}

\newif\ifprofessor
\professorfalse            % <-- troque para \professortrue para incluir o gabarito
\ifdefined\versaoprofessor \professortrue \fi

\usepackage[T1]{fontenc}
\usepackage[utf8]{inputenc}
\usepackage{lmodern}
\usepackage{textcomp}
\usepackage[provide=*,brazilian]{babel}
\usepackage[a4paper,margin=2.5cm]{geometry}
\usepackage{fancyvrb}
\usepackage{upquote}
\usepackage{booktabs}
\usepackage{enumitem}
\usepackage{comment}
\usepackage{url}
\usepackage[hidelinks]{hyperref}

\fvset{xleftmargin=1.5em,fontsize=\small}
\setlist{itemsep=0.2em,topsep=0.4em}
\setlength{\parskip}{0.35em}
\setlength{\emergencystretch}{3em}
\urlstyle{tt}

\newcommand{\arq}[1]{\path{#1}}
\newcommand{\ui}[1]{\textit{#1}}
\newcommand{\tecla}[1]{\textsf{#1}}

\ifprofessor
  \specialcomment{resposta}{\begin{quote}\small\textbf{Gabarito.}\ }{\end{quote}}
\else
  \excludecomment{resposta}
\fi

\title{Árvore de máxima verossimilhança com o IQ-TREE 3\\[0.3em]
\large Prática 5: análise particionada no terminal}
\author{Disciplina Métodos Computacionais\\Universidade Estadual Paulista (UNESP)}
\date{\ifprofessor Versão do professor (com gabarito)\else Versão do aluno\fi\\ Novembro de 2026}

\begin{document}
\maketitle

\paragraph{Origem do material.} Esta apostila adapta o tutorial \emph{Construindo uma árvore
de máxima verossimilhança com o software IQ-TREE3}
(\url{https://labesta.org/blog/2025/maximum_likelihood_tree/}) para o terminal do Ubuntu
(WSL) e do macOS, usando a matriz de Hylodidae da Prática~4 dividida em partições. As mudanças
estão no Apêndice~\ref{ap:mudancas}.

\paragraph{Arquivos.} A pasta \arq{pratica_iqtree} contém a matriz concatenada
\arq{hylodidae_concatenado.fasta} e o arquivo \arq{particoes.txt}, os mesmos que você gerou na
Prática~4. Se preferir, substitua-os pelos seus (pasta \arq{matriz} da Prática~4).

% =====================================================================
\section{Máxima verossimilhança, partições e grupo externo}
% =====================================================================

Na \textbf{máxima verossimilhança}, cada árvore possível recebe uma nota: a probabilidade de
observarmos o nosso alinhamento se aquela árvore, com aqueles comprimentos de ramo, fosse a
verdadeira. A árvore escolhida é a de maior verossimilhança. Para calcular essa probabilidade, é
preciso um \textbf{modelo de substituição de nucleotídeos}, que descreve como as bases mudam ao
longo do tempo. O IQ-TREE 3 \cite{wong2026} escolhe o modelo automaticamente com o ModelFinder
\cite{kalyaanamoorthy2017} e, logo em seguida, procura a melhor árvore. O apoio a cada clado é
medido pelo \emph{bootstrap} ultrarrápido \cite{hoang2018}.

Nossos dois marcadores não evoluem do mesmo jeito, por isso vamos dividir a matriz em
\textbf{partições}, cada uma com o seu próprio modelo:
\begin{itemize}
\item o 16S, um gene de RNA ribossomal;
\item a 1ª, a 2ª e a 3ª posições dos códons do COI, separadas porque mudanças na 3ª posição
muitas vezes não alteram o aminoácido e se acumulam bem mais depressa que nas outras.
\end{itemize}
Separar as posições só é possível porque o alinhamento do COI começa na posição~1 do códon,
como você conferiu no MEGA na Prática~4.

Por fim, a árvore de máxima verossimilhança não tem raiz. Para enraizá-la, indicamos um
\textbf{grupo externo}: \emph{Limnomedusa macroglossa}, que não pertence a Hylodidae.

% =====================================================================
\section{Instalação}
% =====================================================================

Vamos baixar o IQ-TREE 3 (versão 3.1.4) direto do GitHub pelo terminal e descompactá-lo na
sua pasta pessoal (\verb|~|). Escolha os comandos do seu sistema.

\paragraph{Windows (Ubuntu no WSL).}
\begin{Verbatim}[fontsize=\footnotesize,xleftmargin=0pt]
$ cd ~
$ curl -LO \
  https://github.com/iqtree/iqtree3/releases/download/v3.1.4/iqtree-3.1.4-Linux-intel.tar.gz
$ tar -xzf iqtree-3.1.4-Linux-intel.tar.gz
\end{Verbatim}

\paragraph{macOS.}
\begin{Verbatim}[fontsize=\footnotesize,xleftmargin=0pt]
$ cd ~
$ curl -LO \
  https://github.com/iqtree/iqtree3/releases/download/v3.1.4/iqtree-3.1.4-macOS.zip
$ unzip iqtree-3.1.4-macOS.zip
\end{Verbatim}

O \verb|curl -LO| baixa o arquivo do endereço indicado, e o \verb|tar -xzf| (ou o
\verb|unzip|) o descompacta. Agora existe, na sua pasta pessoal, a pasta
\arq{iqtree-3.1.4-Linux-intel} (ou \arq{iqtree-3.1.4-macOS}), com o programa dentro de
\arq{bin}.

% =====================================================================
\section{Preparando a pasta de trabalho}
% =====================================================================

Como no tutorial original, o programa e os dados ficarão na mesma pasta. Vá até a pasta
\arq{pratica_iqtree} (veja a Prática~4: no Windows, ela fica dentro de \arq{/mnt/c/Users/...};
no macOS, digite \verb|cd|, um espaço, e arraste a pasta para o Terminal) e copie o programa
para ela. O ponto no fim do comando significa ``para a pasta atual'':
\begin{Verbatim}
$ cp ~/iqtree-3.1.4-Linux-intel/bin/iqtree3 .
\end{Verbatim}
No macOS, o caminho é \verb|~/iqtree-3.1.4-macOS/bin/iqtree3|. Confira:
\begin{Verbatim}
$ ./iqtree3 --version
\end{Verbatim}
A primeira linha da resposta deve mostrar a versão: \verb|IQ-TREE version 3.1.4|. O \verb|./|
diz ao terminal que o programa está na pasta atual.

% =====================================================================
\section{O arquivo de partições}
% =====================================================================

Abra o \arq{particoes.txt} da Prática~4 no Sublime Text ou no Notepad++. Ele diz quais
colunas da matriz são do 16S e quais são do COI. Com o MAFFT v7.471, por exemplo:
\begin{Verbatim}
DNA, 16S = 1-2475
DNA, COI = 2476-3187
\end{Verbatim}
Crie no editor um arquivo novo, escreva o texto abaixo e salve-o na pasta
\arq{pratica_iqtree} com o nome \arq{particoes.nex}:
\begin{Verbatim}
#nexus
begin sets;
  charset p16S     = 1-2475;
  charset COI_pos1 = 2476-3187\3;
  charset COI_pos2 = 2477-3187\3;
  charset COI_pos3 = 2478-3187\3;
end;
\end{Verbatim}
Cada \verb|charset| dá um nome a um conjunto de colunas. O \verb|\3| significa ``de 3 em 3'':
\verb|2476-3187\3| são as colunas 2476, 2479, 2482\ldots, ou seja, todas as 1ª posições dos
códons do COI; começando em 2477 temos as 2ª posições e, em 2478, as 3ª. O nome \verb|p16S|
começa com uma letra porque, no formato NEXUS, um nome não pode começar com número.

\textbf{Se os números do seu \arq{particoes.txt} forem outros}, use os seus: o 16S vai de 1
até o último número dele, e as três partições do COI começam no primeiro número do COI, no
seguinte e no outro, todas terminando no último número do COI. Não esqueça os pontos e
vírgulas.

% =====================================================================
\section{Rodando a análise}
% =====================================================================

Na pasta \arq{pratica_iqtree}, digite:
\begin{Verbatim}
$ ./iqtree3 -s hylodidae_concatenado.fasta -p particoes.nex \
      -o Limnomedusa_macroglossa -B 1000 --prefix hylodidae
\end{Verbatim}
Como na Prática~4, a barra invertida indica que o comando continua na linha de baixo. Por
partes:
\begin{itemize}
\item \verb|-s hylodidae_concatenado.fasta|: o alinhamento;
\item \verb|-p particoes.nex|: as partições. O IQ-TREE escolhe um modelo para cada uma e
estima uma única árvore para todas;
\item \verb|-o Limnomedusa_macroglossa|: o grupo externo, escrito exatamente como na matriz;
\item \verb|-B 1000|: \emph{bootstrap} ultrarrápido com 1000 réplicas;
\item \verb|--prefix hylodidae|: o começo do nome de todos os arquivos de resultado.
\end{itemize}
A análise leva poucos segundos (cerca de 12~s em um computador comum). Enquanto ela roda,
aparecem várias mensagens; algumas começam com \verb|WARNING| (por exemplo, sobre sequências
com muitos dados faltantes) e, nesta prática, podem ser ignoradas. No Ubuntu, também aparece no
início a mensagem \verb|libarcher.so: cannot open shared object file|, que não atrapalha
nada. A análise terminou quando
aparecer a lista de arquivos de resultado e o \emph{prompt} voltar.

% =====================================================================
\section{Os resultados}
% =====================================================================

O IQ-TREE cria vários arquivos começando com \arq{hylodidae}. Os três principais são:
\begin{itemize}
\item \arq{hylodidae.treefile}: a árvore de máxima verossimilhança, com os valores de
\emph{bootstrap} nos nós, no formato Newick;
\item \arq{hylodidae.iqtree}: o relatório da análise, em texto. Na seção \verb|ModelFinder|,
a linha \verb|Best-fit model according to BIC| traz o modelo escolhido para cada partição;
\item \arq{hylodidae.log}: tudo o que apareceu na tela durante a análise.
\end{itemize}

\paragraph{Vendo a árvore.} Use o FigTree (\url{https://github.com/rambaut/figtree/releases};
no Windows, ele precisa do Java instalado). No Windows, o FigTree é um programa comum, fora do
Ubuntu, e enxerga a pasta \arq{pratica_iqtree} normalmente.
\begin{enumerate}
\item Abra o arquivo \arq{hylodidae.treefile} (\ui{File > Open}). Quando o FigTree pedir um
nome para os rótulos dos nós, digite \verb|bootstrap| e clique em \ui{OK}.
\item Clique no ramo de \emph{Limnomedusa macroglossa} e no botão \ui{Reroot}, para que o
grupo externo fique de um lado da raiz e Hylodidae do outro.
\item No painel da esquerda, marque \ui{Node Labels} e, em \ui{Display}, escolha
\ui{bootstrap}.
\end{enumerate}
Compare a sua árvore com a de um colega: os valores de \emph{bootstrap} podem diferir um
pouco, porque as réplicas são sorteadas.

% =====================================================================
\section*{Perguntas}
\addcontentsline{toc}{section}{Perguntas}
% =====================================================================

\begin{enumerate}
\item Qual modelo foi escolhido para cada uma das quatro partições?
\item Os gêneros \emph{Crossodactylus}, \emph{Hylodes} e \emph{Phantasmarana} saíram
monofiléticos? Com que valor de \emph{bootstrap}?
\end{enumerate}

\begin{resposta}
1. Pelo BIC: \verb|p16S| GTR+F+I+G4; \verb|COI_pos1| TN+I+G4; \verb|COI_pos2| F81+F;
\verb|COI_pos3| TN+F+G4. A 2ª posição do COI não tem nenhum sítio variável (o IQ-TREE avisa:
\verb|No variant sites in partition COI_pos2|), por isso recebe o modelo mais simples.
2. Sim, os três. \emph{Crossodactylus}: 100; \emph{Hylodes}: 85--87; \emph{Phantasmarana}:
99--100 (faixas obtidas em seis execuções). \emph{Megaelosia goeldii}, a única espécie do
gênero, aparece como grupo-irmão de \emph{Phantasmarana} (70--77), e
\emph{Hylodes}~+~\emph{Megaelosia}~+~\emph{Phantasmarana} tem 99--100.
\end{resposta}

% =====================================================================
\appendix
\section{Problemas comuns}\label{ap:problemas}
% =====================================================================

\begin{description}[style=nextline,leftmargin=1.5em]
\item[\texttt{iqtree3: command not found}] Faltou o \verb|./| antes do nome, ou o programa
não foi copiado para a pasta atual (confira com \verb|ls|).
\item[macOS: o sistema diz que não pode verificar o desenvolvedor] Acontece quando o
programa é baixado pelo navegador, e não pelo \verb|curl|. Rode
\verb|xattr -d com.apple.quarantine iqtree3| na pasta \arq{pratica_iqtree}, ou libere o
programa em \ui{Ajustes do Sistema > Privacidade e Segurança}.
\item[\texttt{ERROR: Too large site ID}] Algum número do \arq{particoes.nex} é maior que o
número de colunas da matriz. Confira os números com o seu \arq{particoes.txt}.
\item[\texttt{ERROR: Alignment does not have specified outgroup taxon}] O nome depois do
\verb|-o| está diferente do nome na matriz.
\item[O IQ-TREE não encontra \texttt{particoes.nex}] No Windows, o editor pode ter salvo o
arquivo como \arq{particoes.nex.txt}. Confira com \verb|ls| e, se for o caso, renomeie com
\verb|mv particoes.nex.txt particoes.nex|.
\item[\texttt{Use -redo option if you really want to redo the analysis}] Já existe uma
análise com o mesmo \verb|--prefix|. Acrescente \verb|-redo| ao comando para refazê-la, ou use
outro prefixo.
\end{description}

% =====================================================================
\section{Alterações em relação à versão anterior}\label{ap:mudancas}
% =====================================================================

\begin{enumerate}
\item \textbf{Terminal.} O \arq{iqtree3-click.exe} do Windows, em que as opções eram digitadas
uma a uma (\verb|-s|, \verb|e|, o arquivo, \verb|e|, \verb|-B|\ldots), foi substituído por um
único comando no terminal do Ubuntu (WSL) ou do macOS, com o IQ-TREE 3.1.4 baixado pelo
próprio terminal.
\item \textbf{Dados.} No lugar do 16S de \emph{Pristimantis rupicola}, é usada a matriz
concatenada de Hylodidae (16S + COI) da Prática~4.
\item \textbf{Partições e grupo externo.} Incluídos o arquivo de partições (16S e as três
posições do códon do COI) e o grupo externo (\emph{Limnomedusa macroglossa}).
\item \textbf{Revisão.} Correção do endereço do IQ-TREE (``\texttt{https://https://}'') e da
menção ao IQ-TREE2; explicação dos arquivos de resultado e duas perguntas.
\end{enumerate}

\begin{thebibliography}{9}
\bibitem{hoang2018} Hoang, D.~T., Chernomor, O., von Haeseler, A., Minh, B.~Q. \& Vinh,
L.~S. 2018. UFBoot2: improving the ultrafast bootstrap approximation. \emph{Molecular Biology
and Evolution} 35(2): 518--522.
\bibitem{kalyaanamoorthy2017} Kalyaanamoorthy, S., Minh, B.~Q., Wong, T.~K.~F., von Haeseler,
A. \& Jermiin, L.~S. 2017. ModelFinder: fast model selection for accurate phylogenetic
estimates. \emph{Nature Methods} 14: 587--589.
\bibitem{wong2026} Wong, T.~K.~F., Ly-Trong, N., Ren, H., Demotte, P., Banos, H., Roger,
A.~J., Susko, E., Bielow, C., De~Maio, N., Goldman, N., Hahn, M.~W., dos~Reis, M., Vinh,
L.~S., Huttley, G., Lanfear, R. \& Minh, B.~Q. 2026. IQ-TREE 3: phylogenomic inference
software using complex evolutionary models. \emph{Molecular Biology and Evolution}, msag117.
\end{thebibliography}

\end{document}
